\paragraph{A controlled continuum trajectory in the same matter action.}
The charged action contains a nonlinear real-matter sector for which a
spatial solution estimate can be proved. The electric and magnetic fields
vanish in this sector; the scalar quartic interaction remains active.
The construction first specifies a uniformly regular refinement of the
actual solid, then compares its trajectories with a supplied smooth
continuum solution. It does not assume that consistency on smooth test
functions controls arbitrary finite-element errors.

\begin{proposition}[A uniform conforming cone refinement]
\label{prop:whitney-uniform-cone-refinement}
There is a nested conforming refinement \(\mathcal T_n\),
\(n=2^k\), of the twenty-tetrahedron cone with \(20n^3\) cells.
It respects every macro-face and has
\[
 \frac{1}{2n}<\operatorname{diam}K<\frac6n,
 \qquad r_K>\frac{1}{12n},
 \qquad \frac{\operatorname{diam}K}{r_K}<72
 \quad(K\in\mathcal T_n).
\]
Here \(r_K\) is the inradius. The constants are independent of \(n\).
\end{proposition}
\begin{proof}
Give the macro-vertices one global order. On an ordered tetrahedron
\((v_0,v_1,v_2,v_3)\), use cumulative barycentric coordinates
\[
 x=v_0+B y,\qquad
 B=(v_1-v_0,\ v_2-v_1,\ v_3-v_2),\qquad
 0\leq y_3\leq y_2\leq y_1\leq1.
\]
Intersect this ordered simplex with the Kuhn triangulation of the
\(1/n\) cubical lattice. In a cube based at \(b/n\), its tetrahedra
have vertices
\[
 \frac bn,\quad\frac{b+e_{\pi(1)}}n,\quad
 \frac{b+e_{\pi(1)}+e_{\pi(2)}}n,\quad
 \frac{b+(1,1,1)}n,
 \qquad \pi\in S_3.
\]
Retain those contained in the ordered simplex. The hyperplanes
\(y_i=j/n\) and \(y_i-y_l=j/n\) describe this triangulation.
They include the ordered-simplex boundary and subdivide its interior into
disjoint open tetrahedra. Their volume is \(1/(6n^3)\), giving
\(n^3\) cells per macro-tetrahedron. The hyperplanes at scale \(n\)
occur at scale \(2n\), so these refinements are nested.
On a macro-face a cumulative coordinate becomes an endpoint or two
adjacent coordinates agree. The induced arrangement is precisely the
lower-dimensional construction in the inherited vertex order. Thus both
incident macro-cells give the same face subdivision. This is the standard
edgewise construction \cite{EdelsbrunnerGrayson2000}; the coordinate
description makes its compatibility explicit here.

Every cell is a translate of \(BP\widehat K/n\), where \(P\)
is a coordinate permutation and \(\widehat K\) is one fixed Kuhn
tetrahedron. Its diameter is \(\sqrt3\) and inradius is
\(1/(2+2\sqrt2)\). The centre is the first vertex of every ordered
macro-cell. With \(\varphi=(1+\sqrt5)/2\), its matrix satisfies
\(\|B\|_F^2=\varphi+10<12\) and
\(|\det B|=3+\sqrt5>5\). Hence
\(\sigma_{\max}(B)<\sqrt{12}\) and
\(\sigma_{\min}(B)>5/12\): use
\(|\det B|\leq\sigma_{\min}(B)\sigma_{\max}(B)^2\).
Linear maps bound diameters above and below by these singular values,
and map an inscribed ball to an ellipsoid containing a ball of radius
\(\sigma_{\min}(B)r_{\widehat K}\). The displayed estimates follow.
\end{proof}

\begin{lemma}[The real sector is invariant under the full action]
\label{lem:whitney-real-sector-invariant}
On any conforming tetrahedral refinement, set \(a=\phi=0\) and
take real scalar coefficients and real velocities. Solutions of
\begin{equation}
 (\ddot u_n,v_n)+(\nabla u_n,\nabla v_n)
       +m^2(u_n,v_n)+g(u_n^3,v_n)=0
 \quad\text{for all }v_n\in V_n
 \label{eq:whitney-real-discrete-wave}
\end{equation}
solve all equations of the original charged action, for any real \(e\).
Here \(V_n\) is the continuous piecewise-affine real space on
\(\mathcal T_n\), including its boundary degrees of freedom.
All Gauss constraints and electromagnetic fields vanish exactly.
\end{lemma}
\begin{proof}
At \(a=0\) the dressed interpolation is exactly \(u_n=W_0u\).
Its variation in any real edge direction is
\(ie\sum_i\lambda_i\beta_i u_i\), which is purely imaginary.
Its time and spatial derivatives are also purely imaginary, including
the terms containing the edge-direction velocity and the real scalar
velocity. The additional term \(-ieB_nu_n\) in the covariant spatial
variation is imaginary. Every scalar kinetic, gradient and potential
first variation pairs these terms with real fields and takes the real
part, giving zero. Variations in imaginary scalar coefficients and in
the scalar potential vanish for the same reason. Maxwell terms vanish
because both field strengths are zero. Real scalar variations give the
displayed wave equation after cancelling the common factor two.
This argument tests all variations of the full action; it does not drop
the dressing derivative. Uniqueness of the finite classical equations
then makes this sector invariant. Its local charge and current vanish,
even when \(g>0\) gives nonlinear scalar motion.
\end{proof}

Let \(\Omega\) be the fixed solid and let
\(a_1(v,w)=(\nabla v,\nabla w)+(v,w)\).
Define the \(H^1\) Ritz projection by
\begin{equation}
 a_1(R_nv,v_n)=a_1(v,v_n)\qquad(v_n\in V_n).
 \label{eq:whitney-real-ritz-projection}
\end{equation}
The added mass term controls the constant mode; no boundary value is
fixed by this definition. The uniform mesh gives
\begin{equation}
 \|R_nv-v\|_{H^1}\leq Cn^{-1}\|v\|_{H^2},
 \qquad \|R_nv\|_{H^1}\leq\|v\|_{H^1}.
 \label{eq:whitney-real-ritz-estimate}
\end{equation}
Indeed, orthogonal projection is a contraction in \(a_1\), and its
error is no larger than the piecewise-affine interpolation error. The
latter follows by scaling the fixed reference tetrahedra and the usual
\(H^2\)-to-\(H^1\) affine interpolation estimate. In three dimensions
\(H^2\) has continuous representatives, so its nodal values are defined.
No \(L^2\) duality estimate or elliptic regularity assertion is needed.

\begin{theorem}[Conditional nonlinear matter trajectory convergence]
\label{thm:whitney-real-continuum-trajectory}
Fix \(T<\infty\), \(m^2>0\) and \(g\geq0\). Suppose a real
continuum solution satisfies
\begin{equation}
 u_{tt}-\Delta u+m^2u+gu^3=0,\qquad
 \partial_\nu u=0\text{ on }\partial\Omega,
 \qquad u\in C^2([0,T];H^2(\Omega)).
 \label{eq:whitney-real-continuum-reference}
\end{equation}
The equation and boundary condition may equivalently be stated weakly
against every \(H^1(\Omega)\) test function. Let \(u_n\) solve
the finite equation with \(u_n(0)=R_nu(0)\) and
\(\dot u_n(0)=R_nu_t(0)\). Then
\begin{equation}
 \sup_{0\leq t\leq T}
 \left(\|u_n(t)-u(t)\|_{H^1}
       +\|\dot u_n(t)-u_t(t)\|_{L^2}\right)
 \leq C_T n^{-1}.
 \label{eq:whitney-real-trajectory-error}
\end{equation}
The constant depends on the fixed domain, couplings, time window and
displayed reference norms, but not on the refinement. These are
trajectories of the full charged action in its invariant real sector.
The theorem assumes the indicated smooth continuum reference exists.
\end{theorem}
\begin{proof}
The finite equation conserves
\[
 E_n=\tfrac12\|\dot u_n\|_2^2+\tfrac12\|\nabla u_n\|_2^2
             +\tfrac{m^2}{2}\|u_n\|_2^2+\tfrac g4\|u_n\|_4^4.
\]
Its initial value is uniformly bounded: use Ritz contraction and the
fixed-domain embeddings \(H^1\hookrightarrow L^4,L^6\).
Positive \(m^2\) therefore bounds \(\|u_n\|_{H^1}\) and
\(\|\dot u_n\|_2\) uniformly on the whole interval. For each finite
mesh the positive mass matrix and this bound also give global classical
existence of the polynomial ordinary differential equation.

Put \(\eta=R_nu-u\), \(\theta=u_n-R_nu\).
The time-independent bounded projection commutes with both time
derivatives. Equation~\eqref{eq:whitney-real-ritz-estimate} bounds
\(\eta,\eta_t,\eta_{tt}\) in \(H^1\) by \(C/n\).
Subtracting the full continuum weak equation gives, for every
\(v_n\in V_n\),
\[
 (\ddot\theta,v_n)+a_1(\theta,v_n)
 =-(\eta_{tt},v_n)-(m^2-1)(\theta+\eta,v_n)
                    -g(u_n^3-u^3,v_n).
\]
The Ritz identity cancelled the spatial derivative of \(\eta\);
there is no inverse estimate on \(\dot\theta\).
The cubic factorization and H\"older's inequality imply
\[
 \|v^3-w^3\|_2
 \leq C(\|v\|_{H^1}^2+\|w\|_{H^1}^2)
                      \|v-w\|_{H^1}.
\]
Apply this with \(v=u_n,w=u\), using the uniform bounds above.
Testing with \(\dot\theta\) and writing
\(\mathcal E=\tfrac12\|\dot\theta\|_2^2+
\tfrac12\|\theta\|_{H^1}^2\) yields
\(\dot{\mathcal E}\leq C\mathcal E+C n^{-2}\).
The chosen initialization gives \(\mathcal E(0)=0\).
Gronwall's inequality and the projection estimates prove the claim.
For other uniformly energy-bounded initializations, the same proof adds
\(C_T(\|u_n(0)-R_nu(0)\|_{H^1}+
\|\dot u_n(0)-R_nu_t(0)\|_2)\) to the right-hand side.
\end{proof}

The smooth-reference class has nonzero localized examples. Choose nonzero
\(C_c^\infty\) real data in a ball strictly inside the cone and first solve
on \(\mathbb R^3\). The standard local semilinear Klein--Gordon theorem
and finite propagation speed \cite{TaoNonlinearWaves} give a smooth solution
on a positive interval. At the regularity needed here, local existence
also follows directly from the Klein--Gordon Duhamel map in
\(C H^4\cap C^1H^3\), since the cubic is locally Lipschitz from
\(H^4\) to \(H^3\); the equation gives \(u_{tt}\in C H^2\).
Take the time window shorter than both this existence interval and the
distance of the initial support to the cone boundary. Finite propagation
then leaves an open zero-field collar at that boundary, so restriction to
\(\Omega\) satisfies the stated homogeneous Neumann condition. This uses
standard local continuum PDE input, not a source-derived physical state.

This result supplies an actual nonlinear spatial trajectory comparison,
with zero electromagnetic current. Electrically active complex matter
requires a further stability argument for its field-dependent kinetic
metric and constraints. The present theorem supplies neither that argument
nor physical geometry, a calibrated clock or a physical matter species.

\path{code/electromagnetism/whitney_real_continuum.py} constructs the
ordered refinements, checks a quadratic Ritz comparator with all boundary
variations free, and integrates a compact real pulse under the autonomous
finite quartic-wave equation. Its exact polynomial cell integrals are
independently reconstructed by the verifier. Finite mesh and numerical
time-step comparisons test the implementation; the uniform geometric and
continuum error statements are the analytic proofs above, not numerical
error certificates or Lean formalizations.
